% Figura: sucesiones que definen la delta de Dirac.
\begin{figure}[ht]
\centering
\pgfplotsset{deltaax/.style={tema, width=0.335\textwidth, height=4.3cm, xmin=-1.3, xmax=1.3, ymin=0, ymax=4.4,
   xtick={-1,0,1}, ytick={0,1,2,4}, xlabel={$x$}, grid=major}}
\begin{tikzpicture}
  \begin{axis}[deltaax, name=A, ylabel={$f_\alpha(x)$}, title={\small\scshape (a)\; pulsos}]
    \addplot[azulclaro!80!azulacero, line width=1.1pt, const plot, mark=none] coordinates {(-1.3,0) (-0.5,0) (-0.5,1) (0.5,1) (0.5,0) (1.3,0)};
    \addplot[azulacero, line width=1.1pt, const plot, mark=none] coordinates {(-1.3,0) (-0.25,0) (-0.25,2) (0.25,2) (0.25,0) (1.3,0)};
    \addplot[azulprof, line width=1.1pt, const plot, mark=none] coordinates {(-1.3,0) (-0.125,0) (-0.125,4) (0.125,4) (0.125,0) (1.3,0)};
  \end{axis}
  \begin{axis}[deltaax, at={(A.east)}, anchor=west, xshift=0.6cm, yticklabels={}, title={\small\scshape (b)\; triángulos}]
    \addplot[azulclaro!80!azulacero, line width=1.1pt] coordinates {(-1.3,0) (-1,0) (0,1) (1,0) (1.3,0)};
    \addplot[azulacero, line width=1.1pt] coordinates {(-1.3,0) (-0.5,0) (0,2) (0.5,0) (1.3,0)};
    \addplot[azulprof, line width=1.1pt] coordinates {(-1.3,0) (-0.25,0) (0,4) (0.25,0) (1.3,0)};
  \end{axis}
  \begin{axis}[deltaax, at={(A.east)}, anchor=west, xshift=5.85cm, yticklabels={}, samples=200, domain=-1.3:1.3,
               title={\small\scshape (c)\; gaussianas}]
    \addplot[azulclaro!80!azulacero, line width=1.1pt] {1/(0.6*sqrt(2*pi))*exp(-x^2/(2*0.6^2))};
    \addplot[azulacero, line width=1.1pt] {1/(0.3*sqrt(2*pi))*exp(-x^2/(2*0.3^2))};
    \addplot[azulprof, line width=1.1pt] {1/(0.1*sqrt(2*pi))*exp(-x^2/(2*0.1^2))};
  \end{axis}
\end{tikzpicture}
\caption{La delta de Dirac como límite ($\alpha\to\infty$) de picos cada vez más estrechos y altos, todos de
\emph{área unidad}. (a)~Pulsos rectangulares de anchura $1/\alpha$ y altura $\alpha$ ($\alpha=1,2,4$).
(b)~Triángulos de base $2/\alpha$ y altura $\alpha$. (c)~Gaussianas $\frac{1}{\sigma\sqrt{2\pi}}\eul^{-x^{2}/2\sigma^{2}}$
con $\sigma=0{,}6;\,0{,}3;\,0{,}1$. Cuando el pico se afila, el área no cambia y el valor en $x=0$ diverge.}
\label{fig:delta}
\end{figure}
